\documentclass[12pt,a4paper]{article}

\usepackage[utf8]{inputenc}
\usepackage[T1]{fontenc}
\usepackage[polish]{babel}
\usepackage{geometry}
\usepackage{array}
\usepackage{longtable}
\usepackage{xcolor}
\usepackage{listings}
\usepackage[hidelinks]{hyperref}

\geometry{margin=2.5cm}
\setlength{\parindent}{0pt}
\setlength{\parskip}{6pt}
\emergencystretch=2em

\lstset{
    basicstyle=\ttfamily\small,
    breaklines=true,
    frame=single,
    columns=fullflexible,
    keywordstyle=\color{blue!60!black},
    commentstyle=\color{gray!70!black}
}

\title{Dokumentacja techniczna projektu\\System zarządzania domową biblioteką}
\author{Sebastian Święs}
\date{czerwiec 2026}

\begin{document}

\maketitle
\tableofcontents
\newpage

\section{Cel i zakres projektu}

Projekt jest konsolową aplikacją w języku C++, która służy do zarządzania domową
biblioteką. Program pozwala przechowywać informacje o autorach, kategoriach oraz
pozycjach bibliotecznych, na przykład książkach. Użytkownik może dodawać, edytować,
usuwać i przeglądać dane, a także wyszukiwać pozycje i sprawdzać proste statystyki.

Aplikacja nie korzysta z bazy danych ani interfejsu graficznego. Dane są zapisywane
w plikach CSV w katalogu \texttt{data}. Dzięki temu projekt pozostaje prosty, a
format danych można łatwo podejrzeć w zwykłym edytorze tekstu.

Dokumentacja opisuje strukturę katalogów, najważniejsze klasy, sposób przechowywania
danych, scenariusze działania programu oraz sposób budowania i uruchamiania projektu.

\section{Technologie}

\begin{longtable}{|p{4cm}|p{10cm}|}
\hline
\textbf{Technologia} & \textbf{Zastosowanie} \\
\hline
C++17 & Implementacja modeli danych, logiki biblioteki oraz menu konsolowego. \\
\hline
STL & Wykorzystanie kontenerów, czyli gotowych struktur danych z biblioteki standardowej, na przykład \texttt{std::vector} i \texttt{std::map}. \\
\hline
CMake & Konfiguracja budowania projektu i tworzenie pliku wykonywalnego \texttt{HomeLibrary}. \\
\hline
CSV & Prosty zapis danych autorów, kategorii i pozycji bibliotecznych. \\
\hline
PowerShell & Pomocniczy skrypt \texttt{run.ps1}, który ułatwia budowanie i uruchamianie aplikacji w VS Code. \\
\hline
\end{longtable}

\section{Struktura katalogów}

Najważniejsze elementy projektu są podzielone według odpowiedzialności:

\begin{lstlisting}
projekt_zaawansowany_cpp/
|-- CMakeLists.txt
|-- run.ps1
|-- data/
|   |-- authors.csv
|   |-- categories.csv
|   `-- items.csv
`-- src/
    |-- main.cpp
    |-- models/
    |   |-- Author.h / Author.cpp
    |   |-- Category.h / Category.cpp
    |   `-- Item.h / Item.cpp
    |-- core/
    |   `-- Library.h / Library.cpp
    `-- utils/
        `-- FileStorage.h / FileStorage.cpp
\end{lstlisting}

Katalog \texttt{models} zawiera klasy opisujące dane. Katalog \texttt{core} zawiera
główną logikę aplikacji, a katalog \texttt{utils} odpowiada za zapis i odczyt plików.
Plik \texttt{main.cpp} zawiera menu tekstowe oraz obsługę danych wpisywanych przez
użytkownika.

\section{Konfiguracja budowania}

Projekt jest konfigurowany przez plik \texttt{CMakeLists.txt}. Wymagany jest standard
C++17. Plik wykonywalny ma nazwę \texttt{HomeLibrary}. Do programu dołączane są
wszystkie pliki źródłowe z katalogów \texttt{models}, \texttt{core} i \texttt{utils}.

Przykładowe polecenia budowania:

\begin{lstlisting}
cmake -B build
cmake --build build
.\build\HomeLibrary.exe
\end{lstlisting}

W projekcie znajduje się też skrypt \texttt{run.ps1}. Jest przydatny szczególnie w
VS Code, ponieważ próbuje użyć CMake, a jeśli to się nie uda, kompiluje projekt przez
\texttt{g++}.

\begin{lstlisting}
powershell -ExecutionPolicy Bypass -File .\run.ps1
\end{lstlisting}

\section{Ogólna architektura}

Program można podzielić na cztery warstwy:

\begin{itemize}
    \item warstwa interfejsu konsolowego w pliku \texttt{main.cpp},
    \item warstwa logiki biblioteki w klasie \texttt{Library},
    \item warstwa modeli danych: \texttt{Author}, \texttt{Category}, \texttt{Item},
    \item warstwa zapisu i odczytu danych w klasie \texttt{FileStorage}.
\end{itemize}

Menu konsolowe pobiera dane od użytkownika i wywołuje metody klasy \texttt{Library}.
Klasa \texttt{Library} przechowuje obiekty w pamięci programu i wykonuje operacje
na listach. Klasa \texttt{FileStorage} zapisuje aktualny stan do plików CSV oraz
wczytuje dane przy starcie programu.

\section{Model danych}

\subsection{Autor}

Klasa \texttt{Author} przechowuje identyfikator, imię i nazwisko autora.
Udostępnia metody dostępowe, metody zmiany pól oraz konwersję do formatu CSV.

\begin{lstlisting}
id;firstName;lastName
\end{lstlisting}

Przykład:

\begin{lstlisting}
3;Olga;Tokarczuk
\end{lstlisting}

\subsection{Kategoria}

Klasa \texttt{Category} opisuje kategorię pozycji bibliotecznej. Każda kategoria
ma identyfikator oraz nazwę.

\begin{lstlisting}
id;name
\end{lstlisting}

Przykład:

\begin{lstlisting}
2;Fantasy
\end{lstlisting}

\subsection{Pozycja biblioteczna}

Klasa \texttt{Item} opisuje pojedynczą pozycję w bibliotece. Zawiera tytuł, odwołanie
do autora i kategorii, ocenę, status oraz opcjonalny opis.

\begin{lstlisting}
id;title;authorId;categoryId;rating;status;description
\end{lstlisting}

Przykład:

\begin{lstlisting}
4;Solaris;5;3;9.300000;przeczytane;Klasyka polskiej fantastyki naukowej.
\end{lstlisting}

Relacja między pozycją, autorem i kategorią jest realizowana przez pola
\texttt{authorId} oraz \texttt{categoryId}. Program nie używa mechanizmu kluczy
obcych z bazy danych, ale podczas dodawania lub edycji pozycji sprawdza, czy wybrany
autor i kategoria istnieją.

\section{Klasa Library}

Klasa \texttt{Library} jest głównym miejscem przechowywania danych w czasie działania
programu. Wewnątrz ma trzy kolekcje:

\begin{lstlisting}
std::vector<Author> authors;
std::vector<Category> categories;
std::vector<Item> items;
\end{lstlisting}

Klasa udostępnia operacje CRUD dla autorów, kategorii i pozycji. Identyfikatory nowych
rekordów są wyznaczane na podstawie największego aktualnego ID, powiększonego o jeden.
Przykładowo metoda \texttt{getNextAuthorId()} przegląda listę autorów i zwraca kolejną
wolną wartość.

W klasie znajdują się też metody wyszukiwania:

\begin{itemize}
    \item wyszukiwanie pozycji po fragmencie tytułu,
    \item pobieranie pozycji konkretnego autora,
    \item pobieranie pozycji z wybranej kategorii,
    \item filtrowanie pozycji po statusie.
\end{itemize}

Statystyki są liczone na podstawie aktualnej listy pozycji. Program wyznacza między
innymi średnią ocenę wszystkich pozycji, liczbę pozycji w kategoriach oraz średnią ocenę
dla konkretnej kategorii.

\section{Obsługa plików CSV}

Za zapis i odczyt danych odpowiada klasa \texttt{FileStorage}. Każdy typ danych ma
oddzielny plik:

\begin{longtable}{|p{4cm}|p{10cm}|}
\hline
\textbf{Plik} & \textbf{Zawartość} \\
\hline
\texttt{authors.csv} & Lista autorów. \\
\hline
\texttt{categories.csv} & Lista kategorii. \\
\hline
\texttt{items.csv} & Lista pozycji bibliotecznych. \\
\hline
\end{longtable}

Przy uruchomieniu aplikacja próbuje wczytać dane z katalogu \texttt{data}. Przy wyjściu
dane są automatycznie zapisywane. W menu dostępne są również ręczne opcje zapisu
i wczytania.

W obecnej wersji programu zapis został zabezpieczony przed prostą sytuacją, w której
pusty stan programu mógłby nadpisać istniejące pliki danych. Jest to przydatne, gdy
program zostanie uruchomiony z innego katalogu roboczego lub gdy nie uda się poprawnie
wczytać danych.

\section{Interfejs konsolowy}

Interfejs użytkownika znajduje się w pliku \texttt{main.cpp}. Jest zbudowany z głównego
menu oraz trzech podmenu:

\begin{itemize}
    \item zarządzanie autorami,
    \item zarządzanie kategoriami,
    \item zarządzanie pozycjami.
\end{itemize}

Dodatkowo w menu głównym znajdują się statystyki oraz ręczny zapis i odczyt danych.
Program korzysta z prostych funkcji pomocniczych do wczytywania wartości liczbowych
i tekstowych. Dzięki temu błędne dane wpisane przez użytkownika nie powinny kończyć
programu, tylko powodować wyświetlenie komunikatu.

Przykładowo ocena pozycji musi być liczbą z zakresu od 0 do 10, a wybór statusu jest
ograniczony do trzech opcji:

\begin{itemize}
    \item \texttt{do przeczytania},
    \item \texttt{w trakcie},
    \item \texttt{przeczytane}.
\end{itemize}

\section{Główne scenariusze działania}

\subsection{Start programu}

\begin{enumerate}
    \item Program ustawia obsługę znaków w konsoli Windows.
    \item Tworzony jest obiekt klasy \texttt{Library}.
    \item Aplikacja szuka katalogu \texttt{data}.
    \item Dane z plików CSV są wczytywane do pamięci.
    \item Wyświetlane jest menu główne.
\end{enumerate}

\subsection{Dodanie pozycji}

\begin{enumerate}
    \item Użytkownik wybiera opcję dodania pozycji.
    \item Program sprawdza, czy istnieje przynajmniej jeden autor i jedna kategoria.
    \item Użytkownik podaje tytuł, autora, kategorię, ocenę, status i opcjonalny opis.
    \item Program waliduje wybrane ID autora i kategorii.
    \item Nowy obiekt \texttt{Item} zostaje dodany do kolekcji w klasie \texttt{Library}.
\end{enumerate}

\subsection{Filtrowanie po statusie}

\begin{enumerate}
    \item Użytkownik wybiera opcję filtrowania po statusie.
    \item Program wyświetla listę dostępnych statusów.
    \item Klasa \texttt{Library} przechodzi po liście pozycji i wybiera te, których status
    zgadza się z wybraną wartością.
    \item Wyniki są wypisywane w konsoli razem z autorem, kategorią i oceną.
\end{enumerate}

\subsection{Zakończenie programu}

\begin{enumerate}
    \item Użytkownik wybiera opcję wyjścia.
    \item Program zapisuje dane do plików CSV.
    \item Po udanym zapisie wyświetlany jest komunikat końcowy.
\end{enumerate}

\section{Walidacja i obsługa błędów}

Program zawiera podstawową walidację danych wpisywanych przez użytkownika. Funkcje
\texttt{readIntInRange()} oraz \texttt{readDoubleInRange()} wymuszają zakres liczb.
Funkcja \texttt{readRequiredLine()} pilnuje, aby wymagane pola tekstowe nie były puste.

Przy edycji pozycji program sprawdza, czy wpisane ID autora lub kategorii jest poprawną
liczbą oraz czy wskazany obiekt istnieje. Jeżeli dane są błędne, poprzednia wartość
zostaje zachowana, a użytkownik dostaje komunikat.

Wczytywanie z CSV jest również zabezpieczone przed prostymi błędami formatu. Jeżeli
wiersz nie ma wymaganej liczby pól albo konwersja liczby się nie uda, tworzony jest
pusty obiekt, który nie jest dodawany do biblioteki.

\section{Ograniczenia projektu}

Projekt jest celowo prosty i konsolowy. Najważniejsze ograniczenia to:

\begin{itemize}
    \item brak graficznego interfejsu użytkownika,
    \item brak prawdziwej relacyjnej bazy danych,
    \item prosty format CSV bez obsługi średników wewnątrz tekstu,
    \item brak rozbudowanych testów automatycznych,
    \item brak kont użytkowników i uprawnień.
\end{itemize}

Te ograniczenia wynikają głównie z zakresu projektu. Dla małej domowej biblioteki
plikowy zapis CSV jest wystarczający i łatwy do zrozumienia.

\section{Podsumowanie}

System zarządzania domową biblioteką jest prostą aplikacją konsolową napisaną w C++17.
Projekt pokazuje wykorzystanie klas, enkapsulacji, kontenerów STL, na przykład \texttt{std::vector}, podziału na moduły
oraz zapisu danych do plików. Najważniejszą klasą logiki jest \texttt{Library}, a za
trwałość danych odpowiada \texttt{FileStorage}. Program realizuje podstawowe operacje
CRUD, wyszukiwanie, filtrowanie po statusie oraz statystyki, dzięki czemu spełnia
założenia małej aplikacji do zarządzania prywatną kolekcją książek.

\end{document}
